\VerdierPhysicalPageAnchor{0069}%
\VerdierSemanticAnchor{ch1-2-2-2-continuation}%
form an $m$-fold complex in $C$.

\VerdierSemanticAnchor{ch1-2-2-2-morphism}
Now let $f\colon Y\longrightarrow Y'$ be a morphism of
$\varphi$-summable $n$-fold complexes in $C$. Set:

\VerdierSemanticAnchor{ch1-2-2-2-eq6}
\begin{equation}
  \left(\int_\varphi f\right)^\tau
  =\bigoplus_{\varphi(\sigma)=\tau}f^\sigma\colon
  \left(\int_\varphi Y\right)^\tau
  \longrightarrow
  \left(\int_\varphi Y'\right)^\tau.
  \tag{2.2.2.6}
\end{equation}

\VerdierSemanticAnchor{ch1-2-2-2-functoriality}
It is clear that the $\left(\int_\varphi f\right)^\tau$ commute with the
differentials of $\int_\varphi Y$ and $\int_\varphi Y'$ and that,
consequently, such a family defines a morphism of $m$-fold complexes
$\int_\varphi f\colon\int_\varphi Y\longrightarrow\int_\varphi Y'$.
It is clear that
$\int_\varphi fg=\int_\varphi f\,\int_\varphi g$ and that
$\int_\varphi\operatorname{id}_Y
=\operatorname{id}_{\int_\varphi Y}$. Consequently $\int_\varphi$ is a
functor from the category of $\varphi$-summable $n$-fold complexes to the
category of $m$-fold complexes. This functor is additive.

\VerdierSemanticAnchor{ch1-2-2-3}
\noindent\textbf{2.2.3.}\quad
Let $\varphi\colon[n]\longrightarrow[m]$ be a map, let $\tau$ be an
element of $\mathbb Z[m]$, and let $i$ be an element of $[m]$. Let $E_1$
(resp. $E_2$) be the sets:

\VerdierSemanticAnchor{ch1-2-2-3-fibers}
\[
  \begin{aligned}
    E_1&=\{\sigma_1\in\mathbb Z[n]\mid\varphi(\sigma_1)=\tau\},\\
    \text{(resp. }E_2
      &=\{\sigma_2\in\mathbb Z[n]\mid
          \varphi(\sigma_2)=\tau+e_i\}\text{)}.
  \end{aligned}
\]
Finally, let $Y$ be an $n$-fold complex in $C$. Set, for
$\sigma_1\in E_1$, $\sigma_2\in E_2$:

\VerdierSemanticAnchor{ch1-2-2-3-eq1}
\begin{equation}
  u_{\sigma_2,\sigma_1}(Y)
  =\begin{cases}
    \varepsilon(\tau,i)\varepsilon(\sigma_1,j)d_Y^{\sigma_1,j},
      &\text{if }\sigma_2-\sigma_1=e_j,\\
    0,&\text{otherwise},
  \end{cases}
  \tag{2.2.3.1}
\end{equation}
whence a matrix of type $E_2\times E_1$:

\VerdierSemanticAnchor{ch1-2-2-3-eq2}
\begin{equation}
  M(\tau,i,Y)
  =\bigl(u_{\sigma_2,\sigma_1}(Y)\bigr)_
    {\sigma_2\in E_2,\,\sigma_1\in E_1}.
  \tag{2.2.3.2}
\end{equation}

\VerdierSemanticAnchor{ch1-2-2-3-matrix-identification}
When $Y$ is $\varphi$-summable, this matrix is the matrix of the morphism
$d_{\int_\varphi Y}^{\tau,i}$.

\VerdierSemanticAnchor{ch1-2-2-3-extensions}
Observe that in general this matrix has only finitely many nonzero entries
in each row and in each column. There are therefore \emph{a priori} at
least two possible extensions of the functor $\int_\varphi$ to arbitrary
complexes when, in $C$, countable products and direct sums
